
This work presents a controlled comparison of uncertainty quantification methods for hallucination detection, finding that UQ method effectiveness is format-dependent. On TruthfulQA multiple-choice format with Llama-3-8B-Instruct:

\begin{itemize}
\item Token-level methods (max probability, choice entropy) achieve AUROC 0.77--0.81
\item Semantic entropy achieves only AUROC 0.5645 despite correct mechanism implementation
\item The gap stems from format mismatch: multiple-choice answers lack semantic content for NLI comparison
\end{itemize}

The practical implication is that method selection should match task format. For multiple-choice tasks, simple confidence-based methods are sufficient and computationally cheaper.

Future work should evaluate semantic entropy on free-form generation tasks (e.g., HaluEval dialogue and summarization) where the format hypothesis predicts improved performance. Cross-model validation and full-scale experiments on the complete TruthfulQA dataset would strengthen the findings.

